\documentclass[11pt,a4paper]{article}
\usepackage[a4paper,margin=16mm]{geometry}
\usepackage[T1]{fontenc}
\usepackage{amsmath,amssymb}
\usepackage{enumitem}

\setlength{\parindent}{0pt}
\setlength{\parskip}{4pt}

\begin{document}

\begin{center}
  {\Large\bfseries Answer Key}\\[2pt]
  {\large Picture Sequences and the \(n\)th Term}
\end{center}

\section*{Part 1: Linear sequences}
\begin{itemize}[leftmargin=*,itemsep=3pt]
\item \textbf{Q1.} \(5n+1\); terms: \(6,11,16,21,26,51\).
\item \textbf{Q2.} \(2n+2\); terms: \(4,6,8,10,12,22\).
\item \textbf{Q3.} (a) \(4n+1\); terms \(5,9,13,17,21,41\). (b) \(2n-1\); terms \(1,3,5,7,9,19\).
\item \textbf{Q4.} (a) \(n=20\). (b) \(n=14\). (c) No; \(n=19.8\) is not a whole number.
\item \textbf{Q5.} (a) \(5,7,9\). (b) \(3,7,11\). (c) \(7,12,17\). Correct drawings matching these counts are acceptable.
\end{itemize}

\section*{Part 2: Square numbers}
\begin{itemize}[leftmargin=*,itemsep=3pt]
\item \textbf{Q6.} (a) \(n^2\); terms \(1,4,9,16,25,100\). (b) \(25=1+3+5+7+9\); the added L-shape has \(2n-1\) dots.
\item \textbf{Q7.} (a) \(n^2+2\); terms \(3,6,11,18,27,102\). (b) \(n(n+1)=n^2+n\); terms \(2,6,12,20,30,110\). (c) \(2n^2\); terms \(2,8,18,32,50,200\). (d) \(n^2-1\); terms \(0,3,8,15,24,99\).
\item \textbf{Q8.} (a) \(1,3,6,10,15,55\). (b) \(n(n+1)\). (c) \(\frac{n(n+1)}{2}\); terms \(1,3,6,10,15,55\).
\item \textbf{Q9.} (a) \(n^2+3\): \(4,7,12\). (b) \(3n^2\): \(3,12,27\). (c) \(n^2+2n=n(n+2)\): \(3,8,15\). Correct drawings matching these counts are acceptable.
\end{itemize}

\section*{Part 3: Pixel art sequences}
\begin{itemize}[leftmargin=*,itemsep=3pt]
\item \textbf{Q10.} \(3n+8\); terms \(11,14,17,20,23,38\). For 50 pixels, \(n=14\). For 100 pixels, impossible.
\item \textbf{Q11.} \(16n^2\); terms \(16,64,144,256,400,1600\). Doubling the pattern number multiplies the pixels by 4.
\item \textbf{Q12.} \(n^2\); terms \(1,4,9,16,25,100\). The pyramid rearranges into an \(n\) by \(n\) square.
\item \textbf{Q13.} Picture: \(n^2\); frame: \(4n+4\); total: \((n+2)^2=n^2+4n+4\). For \(n=1,2,3,4,5,10\): picture \(1,4,9,16,25,100\); frame \(8,12,16,20,24,44\); total \(9,16,25,36,49,144\).
\item \textbf{Q14.} Example: \(n^2+4\); terms \(5,8,13,20,29,104\). Other valid designs are possible.
\end{itemize}

\section*{Part 4: Always, sometimes or never?}
\begin{itemize}[leftmargin=*,itemsep=3pt]
\item \textbf{Q15.} Always.
\item \textbf{Q16.} Never.
\item \textbf{Q17.} Sometimes.
\item \textbf{Q18.} Sometimes.
\item \textbf{Q19.} Sometimes.
\item \textbf{Q20.} Always.
\item \textbf{Q21.} Never.
\item \textbf{Q22.} Never.
\item \textbf{Q23.} Always.
\item \textbf{Q24.} Never.
\item \textbf{Q25.} Always.
\item \textbf{Q26.} Sometimes.
\item \textbf{Q27.} Example: Always \(6n+2\) is even; Sometimes \(3n+2\) is a multiple of 5; Never \(5n+2\) is a multiple of 5.
\end{itemize}

\end{document}